%% 第 26 章：奇异值分解
%% 本文件由 tools/md2tex.py 自动生成，请勿直接编辑。

\chapter{奇异值分解}


\section{为什么需要奇异值分解}

特征值分解要求：

\[
T:V\to V.
\]

而奇异值分解适用于任意有限维内积空间之间的线性映射：

\[
T:V\to W.
\]

它同时在输入空间和输出空间中选择规范正交基，使 $T$ 只沿对应方向进行独立缩放。

\section[从 T*T 出发]{从 $T^*T$ 出发}

算子：

\[
T^*T:V\to V
\]

是自伴正算子。

由谱定理，存在 $V$ 的规范正交基：

\[
(v_1,\ldots,v_n),
\]

使：

\[
T^*Tv_i=\sigma_i^2v_i,
\qquad
\sigma_i\ge0.
\]

其中非负数：

\[
\sigma_i
\]

称为 $T$ 的奇异值。

\begin{jprop}{非零奇异值对应的输出方向}{r:27:1}
设：

\[
\sigma_i>0.
\]

定义：

\[
u_i
=
\frac{T(v_i)}{\sigma_i}.
\]

则：

\[
T(v_i)=\sigma_iu_i.
\]

并且全部这样的 $u_i$ 构成 $W$ 中的规范正交列表。
\end{jprop}

\begin{jproof}{证明思路}
对任意正奇异值对应的 $v_i,v_j$：

\[
\langle T(v_i),T(v_j)\rangle
=
\langle v_i,T^*Tv_j\rangle.
\]

而：

\[
T^*Tv_j=\sigma_j^2v_j.
\]

因为 $v_i$ 彼此规范正交，得到：

\[
\langle T(v_i),T(v_j)\rangle
=
\sigma_i^2\delta_{ij}.
\]

除以对应的 $\sigma_i,\sigma_j$ 后，便得到：

\[
\langle u_i,u_j\rangle=\delta_{ij}.
\]
\end{jproof}

\begin{jprop}{零奇异值与零空间}{r:27:2}
若：

\[
\sigma_i=0,
\]

则：

\[
T^*Tv_i=0.
\]

于是：

\[
\|Tv_i\|^2
=
\langle T^*Tv_i,v_i\rangle
=
0.
\]

因此：

\[
Tv_i=0.
\]

所以，零奇异值对应的右奇异向量恰好张成：

\[
\ker T.
\]

这就是 SVD 中“降维”发生的位置：这些输入方向会被完全压缩掉。
\end{jprop}

\begin{jthm}{奇异值分解定理}{r:27:3}
设：

\[
T:V\to W
\]

是有限维内积空间之间的线性映射，且：

\[
\operatorname{rank}T=r.
\]

则存在 $V$ 的规范正交基：

\[
(v_1,\ldots,v_n),
\]

以及 $W$ 的规范正交基：

\[
(u_1,\ldots,u_m),
\]

使：

\[
T(v_i)=\sigma_i u_i,
\qquad
i=1,\ldots,r,
\]

其中：

\[
\sigma_1,\ldots,\sigma_r>0.
\]

并且：

\[
T(v_i)=0,
\qquad
i=r+1,\ldots,n.
\]

在这两组基下，$T$ 的矩阵为矩形对角矩阵：

\[
\Sigma
=
\begin{pmatrix}
\sigma_1 & 0 & \cdots & 0\\
0 & \sigma_2 & \cdots & 0\\
\vdots & \vdots & \ddots & \vdots\\
0 & 0 & \cdots & \sigma_r\\
0 & 0 & \cdots & 0\\
\vdots & \vdots & & \vdots
\end{pmatrix}.
\]

矩阵形式常写为：

\[
A=U\Sigma V^*.
\]
\end{jthm}

\section{SVD 为什么一定有足够多的基向量}

第一步，谱定理为：

\[
T^*T
\]

提供 $V$ 的完整规范正交特征基。

其中正特征值方向经过 $T$ 后，构造出 $W$ 中规范正交的左奇异向量。

零特征值方向正好组成：

\[
\ker T.
\]

最后，将已有左奇异向量补全为 $W$ 的规范正交基。

因此，SVD 不假设输入输出空间恰好存在配对基（这是我们先前学习的一大误解），它由：

\[
T^*T
\]

的正谱结构主动构造输入基，再由 $T$ 构造输出基。

\section[从 TT* 出发]{从 $TT^*$ 出发}

同样可以从：

\[
TT^*:W\to W
\]

开始。

若：

\[
T^*Tv=\sigma^2v,
\qquad
\sigma>0,
\]

则：

\[
TT^*(Tv)
=
\sigma^2(Tv).
\]

因此，$T^*T$ 与 $TT^*$ 的全部正特征值相同，且重数相同。

它们分别描述同一组奇异模式在输入空间与输出空间中的表现：

\begin{jitem}
\item %
$T^*T$ 给出右奇异向量，即输入方向；
\item %
$TT^*$ 给出左奇异向量，即输出方向。
\end{jitem}

二者在零特征值部分可能维数不同，这正反映了定义域、陪域维数与核、余核结构的差异。
